\documentclass[12pt]{article}
\usepackage[T2A]{fontenc}
\usepackage[utf8]{inputenc}
\usepackage[russian]{babel}
\usepackage{paratype}
\usepackage{amsmath,amssymb}
\usepackage{graphicx}
\graphicspath{{./figures/}}
\usepackage{booktabs}
\usepackage{array}
\usepackage{multirow}
\usepackage{colortbl}
\usepackage{xcolor}
\usepackage{tikz}
\usepackage{geometry}
\usepackage{fancyhdr}
\usepackage{caption}
\usepackage{hyperref}

\geometry{a4paper, margin=2.5cm}

% Define gray for table cells
\definecolor{lightgray}{gray}{0.85}
\definecolor{darkgray}{gray}{0.55}

\pagestyle{fancy}
\fancyhf{}
\fancyhead[L]{\small\thepage\quad Странные информационные структуры генетического кода}
\fancyhead[R]{\small\thepage}
\renewcommand{\headrulewidth}{0.4pt}

\begin{document}

\begin{center}
{\Large\bfseries 2}\\[0.5em]
{\Large\bfseries Являются ли странные информационные структуры\\
генетического кода случайностью\\
или артефактом?}

\vspace{1em}

{\large Александр Д.\ Панов}\\[0.3em]
{\small Научно-исследовательский институт ядерной физики имени Д.В.~Скобельцына Московского государственного университета имени М.В.~Ломоносова}

\vspace{0.5em}

{\large Феликс П.\ Филатов}\\[0.3em]
{\small НИИ вакцин и сывороток имени И.И.~Мечникова,\\
Национальный исследовательский центр эпидемиологии и микробиологии имени Н.Ф.~Гамалеи, Москва, Россия}

\vspace{0.5em}

{\small\itshape Эволюция: экологические, демографические и политические риски 2024, с.\ 23--68\\
DOI: 10.30884/978-5-7057-6399-3\_03}
\end{center}

\vspace{2em}

\noindent\textbf{Аннотация}

\noindent\textit{Универсальный генетический код обладает несколькими странными симметриями, биохимическое значение которых остаётся неясным. Кроме того, из масс аминокислот, кодируемых генетическим кодом, можно построить множество числовых сигнатур, обычно связанных с делимостью на 111. Настоящая статья основана на работе Владимира Щербака и Максима Макукова (2013), в которой впервые было получено большое число таких независимых информационных структур. В данной работе найдено несколько новых информационных структур и показано, что вся их совокупность естественным образом разделяется на пять <<информационных уровней>>, первые два из которых особенно просты и напоминают <<сигналы привлечения внимания>>, известные в практике поиска сигналов внеземного разума в проблеме SETI. Метод анализа, использованный в статье, существенно отличается от подхода Щербака -- Макукова и основан на систематическом использовании <<критерия простоты>>. В статье подробно представлены как новые, так и известные информационные сигнатуры генетического кода и обсуждаются перспективы решения вопроса о случайной или искусственной природе этих странных информационных структур.}

\medskip
\noindent\textit{\textbf{Ключевые слова:} происхождение жизни, астробиология, экзобиология, генетический код.}

\setcounter{page}{23}

\section{Введение}

Заголовок статьи следует понимать буквально: содержание статьи представляет собой подробное объяснение вопроса, поставленного в заголовке. Определённого ответа на этот вопрос не даётся, но рассматриваются возможные сценарии развития темы. В работе устанавливается существование альтернативы <<случайность/искусственность>> в отношении странных информационных структур, которые действительно существуют в генетическом коде, и обосновывается правомерность вопроса, поставленного в заголовке статьи. Искусственность упомянутых информационных структур (вторая из возможностей, представленных в альтернативе) означает внеземное происхождение земной жизни, что переводит всю проблему в область астробиологии и задач SETI.

Идея о том, что генетический код может содержать информацию, сигнализирующую об искусственном вмешательстве в происхождение земной жизни, была впервые сформулирована Г.~Марксом (1979). Он писал, что генетический код благодаря своей чрезвычайной стабильности и консерватизму является наиболее надёжным носителем информации биологической природы, хотя объём информации, который он может содержать, невелик (короче текста из 40 букв в обычном алфавите). Маркс, однако, отмечал, что на момент написания его работы в этом носителе не было обнаружено никакой информации, кроме правил кодирования. Серьёзные поиски такой информации в структурах генетического кода были предприняты в работах Владимира Щербака и Максима Макукова (Щербак и Макуков 2013; Макуков и Щербак 2018), и эти поиски привели к обнаружению большого числа нетривиальных информационных закономерностей. Эти статьи послужили отправной точкой настоящей работы.

Как могла бы выглядеть сигнатура искусственности в генетическом коде? Это могли бы быть числовые соотношения, связанные со структурой кода, которые не выглядят случайными, но чьё естественное происхождение не имеет объяснения; неожиданные симметрии кода и~\textit{т.\,д}. Поиск таких <<особенностей>> кода может вызвать упрёк в нумерологии. Однако как ещё могли бы выглядеть такие информационные сигнатуры? Здесь важно не впасть в невольную подгонку и манипуляцию числами, а использовать соотношения, <<лежащие на поверхности>>, избегать слишком сложных и искусственных комбинаторных конструкций и обращать внимание на то, что в определённом смысле удовлетворяет критерию <<естественности и простоты>>. По крайней мере, с этого следует начинать.

Статьи Щербака и Макукова, однако, начинаются сразу с довольно сложных конструкций. Со своей стороны мы старались строго следовать вышеупомянутому критерию естественности и простоты в явном виде. Наш подход также позволил нам заметить новую особенность необычных информационных закономерностей, обнаруженных в генетическом коде: они естественным образом располагаются на нескольких информационных уровнях в порядке возрастания сложности и в порядке привлечения новых ресурсов для представления информации. Более того, два первых, простейших уровня напоминают <<сигнал привлечения внимания>>, наличие которого часто предполагается в сообщениях внеземного разума и который действительно присутствует в реальных земных сообщениях, уже отправленных внеземному разуму от нас (Зайцев 2008). Наш подход также обеспечил дополнительную поддержку более сложного подхода к теме, который был использован в работе Щербака и Макукова.

В настоящей работе мы последовательно представим все известные нам нетривиальные информационные закономерности генетического кода. Для краткости мы будем просто называть такие структуры сигнатурами, не подразумевая, что они обязательно что-либо означают. Не все конструкции в работе Щербака и Макукова представляются нам безупречными. В нашем обзоре мы представим некоторые новые сигнатуры, а также сигнатуры, найденные самими Щербаком и Макуковым (частично переосмысленные), которые не вызывают возражений, а также сигнатуры, известные из ещё более ранних работ (связанные с так называемой симметрией Румера, см. ниже). По ходу обзора мы будем уточнять, является ли рассматриваемая сигнатура новой или уже известной из предыдущих исследований. Основной порядок изложения будет следовать возрастающей сложности информационных уровней, упомянутых выше.

Наше изложение завершится обсуждением полученных результатов и дальнейших перспектив.

Мы начнём с краткого представления структуры генетического кода, которое также позволит нам ввести необходимую терминологию.

\section{Таблица генетического кода}

Таблица генетического кода, известная из всех учебников (см. рис.~1), была предложена Фрэнсисом Криком (1968). Поскольку каждый кодирующий продукт (аминокислота или сигнал \textbf{стоп}) определяется тройкой нуклеотидов (оснований), называемой также кодоном, таблица кода имеет форму трёхмерной матрицы $4 \times 4 \times 4$, индексированной по каждой стороне четырьмя нуклеотидами, например, в порядке тимин--цитозин--аденин--гуанин (T-C-A-G в однобуквенных обозначениях, см. рис.~2). Каждой тройке нуклеотидов соответствует одна ячейка матрицы, в которой находится аминокислота, кодируемая данным кодоном, или сигнал \textbf{стоп}. Обычная двумерная таблица или матрица имеет строки и столбцы, расположенные в горизонтальной плоскости <<листа бумаги>>; в трёхмерной матрице добавляются вертикальные столбцы (см. рис.~2). Можно сказать, что трёхмерная матрица~--- это двумерная матрица с вертикальными столбцами в её ячейках. Трёхмерную матрицу неудобно представлять на двумерном листе бумаги. Поэтому на рис.~1 третье измерение, соответствующее третьему основанию в кодоне, размещено не в виде вертикального столбца под ячейкой двумерной матрицы, как на рис.~2, а <<лёжа>>, так что каждая из ячеек двумерной таблицы, индексированная первыми двумя основаниями кодона, сама является столбцом из четырёх ячеек, пронумерованных по третьему основанию. Поскольку содержимое этих больших ячеек фактически берётся из вертикальных столбцов трёхмерной матрицы, мы по-прежнему будем называть их столбцами. Таким образом, каждый столбец содержит четыре кодирующих продукта при постоянных первых двух основаниях кодона и меняющемся третьем основании.

\begin{figure}[htbp]
\centering
\includegraphics[width=0.75\textwidth]{fig1_genetic_code_table.png}\\[0.5em]
\includegraphics[width=0.75\textwidth]{fig1_calligram.png}

\caption{Таблица универсального генетического кода в форме Крика и её каллиграмма. Кодирующие продукты (аминокислоты) обозначены однобуквенными символами, триплеты, терминирующие трансляцию, обозначены словом \textbf{стоп}. Серые ячейки относятся к октету~\textbf{I} Румера, остальные~--- к октету~\textbf{II} (см. текст)}
\label{fig:genetic_code}
\end{figure}

Основания T, C, A, G представлены соединениями двух типов: T и C~--- пиримидины, A и G~--- пурины. Мы будем использовать стандартное обозначение Y для пиримидинов и R для пуринов. Далее мы будем обозначать множество оснований (T, C, A) через H, а полное множество оснований (T, C, A, G) через V (см. таблицу~1).

В таблице генетического кода на рис.~1 для обозначения аминокислот используются стандартные однобуквенные обозначения \textbf{F}, \textbf{L},~\ldots\ В таблице~2 приведены полные названия всех 20 аминокислот генетического кода вместе с их стандартными однобуквенными и трёхбуквенными обозначениями. В таблице также приведена полная атомная масса $M$ каждой свободной аминокислоты (иначе называемая массовым числом): сумма атомных масс составляющих её атомов (для каждого атома~--- число протонов вместе с нейтронами в ядре), причём для каждого атома берётся атомная масса его наиболее распространённого стабильного изотопа: для водорода 1, для углерода 12, для азота 14, для кислорода 16 и так далее.

\begin{figure}[htbp]
\centering
\includegraphics[width=0.5\textwidth]{fig2_3d_cube_top.png}\\
\includegraphics[width=0.5\textwidth]{fig2_3d_cube_bottom.png}
\caption{Трёхмерная матрица, индексированная основаниями T, C, A, G}
\label{fig:3d_matrix}
\end{figure}

Каждая аминокислота состоит из постоянной части и боковой цепи (см. рис.~3). Постоянные части всех аминокислот, за единственным исключением пролина (\textbf{P}, \textbf{Pro}), одинаковы и имеют атомную массу 74. Постоянная часть пролина имеет атомную массу 73. Боковые цепи имеют разнообразные атомные массы $M_{Side}$, которые также приведены в таблице~2.

\begin{table}[htbp]
\centering
\caption{Обозначения групп оснований}
\label{tab:notations}
\begin{tabular}{|c|l|}
\hline
\textbf{Обозначение} & \textbf{Группа оснований} \\
\hline
Y & (T, C) -- пиримидины \\
R & (A, G) -- пурины \\
H & (T, C, A) \\
N & (T, C, A, G) -- любое основание \\
M & (T, G) -- первая пара Румера \\
K & (C, A) -- вторая пара Румера \\
\hline
\end{tabular}
\end{table}

Как видно из таблицы генетического кода на рис.~1, в большинстве случаев одному кодирующему продукту соответствуют несколько различных кодонов. Это так называемая вырожденность генетического кода. В некоторых случаях все кодоны одного столбца кодируют один и тот же продукт. Такие столбцы мы будем называть однородными. На рис.~1 они помечены серым цветом. В других случаях один столбец кодирует разные продукты; такие столбцы мы будем называть неоднородными, они оставлены незакрашенными. Мы будем сопоставлять таблице генетического кода простую картинку, которую будем называть каллиграммой. Каллиграмма~--- это квадрат $4 \times 4$, ячейки которого закрашены чёрным, если они соответствуют однородным столбцам таблицы генетического кода (серым на рис.~1), и оставлены незакрашенными, если они соответствуют неоднородным столбцам (незакрашенным на рис.~1). Каллиграмма, соответствующая основной таблице генетического кода, показана на рис.~1 справа. Далее мы будем рассматривать другие представления таблицы генетического кода, которым будут соответствовать другие каллиграммы. Теперь перейдём шаг за шагом к представлению нетривиальных информационных структур генетического кода.

\begin{figure}[htbp]
\centering
\includegraphics[width=0.4\textwidth]{fig3_amino_acid.png}
\caption{Структура аминокислоты: постоянная часть и боковая цепь}
\label{fig:amino_acid}
\end{figure}

\begin{table}[htbp]
\centering
\caption{Аминокислоты и их атомные массы}
\label{tab:amino_acids}
\begin{tabular}{|l|c|c|c|c||l|c|c|c|c|}
\hline
\textbf{Аминокислота} & & & $M_{Side}$ & $M$ & \textbf{Аминокислота} & & & $M_{Side}$ & $M$ \\
\hline
Глицин & \textbf{G} & \textbf{Gly} & 1 & 75 & Аспарагиновая кислота & \textbf{D} & \textbf{Asp} & 59 & 133 \\
Аланин & \textbf{A} & \textbf{Ala} & 15 & 89 & Глутамин & \textbf{Q} & \textbf{Gln} & 72 & 146 \\
Серин & \textbf{S} & \textbf{Ser} & 31 & 105 & Лизин & \textbf{K} & \textbf{Lys} & 72 & 146 \\
Пролин & \textbf{P} & \textbf{Pro} & 42 & 115 & Глутаминовая кислота & \textbf{E} & \textbf{Glu} & 73 & 147 \\
Валин & \textbf{V} & \textbf{Val} & 43 & 117 & Метионин & \textbf{M} & \textbf{Met} & 75 & 149 \\
Треонин & \textbf{T} & \textbf{Tre} & 45 & 119 & Гистидин & \textbf{H} & \textbf{His} & 81 & 155 \\
Цистеин & \textbf{C} & \textbf{Cys} & 47 & 121 & Фенилаланин & \textbf{F} & \textbf{Phe} & 91 & 165 \\
Лейцин & \textbf{L} & \textbf{Leu} & 57 & 131 & Аргинин & \textbf{R} & \textbf{Arg} & 100 & 174 \\
Изолейцин & \textbf{I} & \textbf{Ile} & 57 & 131 & Тирозин & \textbf{Y} & \textbf{Tyr} & 107 & 181 \\
Аспарагин & \textbf{N} & \textbf{Asn} & 58 & 132 & Триптофан & \textbf{W} & \textbf{Trp} & 130 & 204 \\
\hline
\end{tabular}
\end{table}


\section{Первый информационный уровень. Симметрия Румера и связанные с ней симметрии генетического кода}

Не вдаваясь в рассуждения о происхождении кода, Фрэнсис Крик назвал его <<замороженной случайностью>>. Позднее были предложены три различные теории, в определённой степени объясняющие структуру кода (Никитин 2016): теория оптимизации для минимизации ошибок синтеза белков, теория структурного соответствия аминокислот кодонам (ключ--замок) и теория коэволюции кодонов и путей биосинтеза аминокислот. Однако оказалось, что генетический код обладает некоторыми неожиданными формальными свойствами, которые, по общему пониманию, не могут быть просто и однозначно объяснены в рамках современных представлений, связанных с биохимией или эволюцией.

Первым обратил на них внимание Юрий Румер (Румер 2013; Конопельченко и Румер 1975). Пытаясь найти рациональную основу структуры генетического кода, он объединил однородные столбцы в одно множество, а неоднородные~--- в другое. Число столбцов в каждом множестве оказалось равным восьми (см. рис.~1), поэтому эти два множества были названы октетами; мы будем обозначать их римскими цифрами \textbf{I} и \textbf{II}. Октет \textbf{I} соответствует однородным столбцам, октет \textbf{II}~--- неоднородным. Классификация Румера нетривиальна, поскольку соотношение числа столбцов в двух множествах вполне могло бы быть другим и действительно является другим в некоторых существующих редких альтернативных генетических кодах (см. раздел~9.3.).

Совершенно неожиданно оказывается, что октеты \textbf{I} и \textbf{II} связаны простым преобразованием симметрии: $R = (\text{T} \leftrightarrow \text{G}, \text{C} \leftrightarrow \text{A})$, которое мы будем называть преобразованием Румера. При такой перестановке оснований каждый столбец октета \textbf{I} переходит в некоторый столбец октета \textbf{II} и наоборот. Например, столбец CT октета \textbf{I} переходит в столбец AG октета \textbf{II}, столбец AG переходит в CT и так далее.

Наличие этой странной симметрии, с одной стороны, легко обнаруживается, как и само существование октетов, но, с другой стороны, не может быть понято с точки зрения особой роли этой симметрии в функционировании генетического кода. Если предположить, что при наличии двух октетов остальная часть генетического кода устроена случайным образом, то условная вероятность появления симметрии Румера равна 1/256, \textit{т.\,е.}\ эта симметрия выглядит довольно маловероятной, хрупкой и необязательной в функциональном смысле (что подтверждается её отсутствием в некоторых альтернативных генетических кодах).

Следует отметить, что эта симметрия была независимо переоткрыта по крайней мере дважды после Румера (Данквертс и Нойберт 1975; Вильгельм и Николаева 2004), потому что она легко обнаруживается, буквально <<лежит на поверхности>>, но не привлекает внимания биохимиков, поскольку связь между этой симметрией и функциями генетического кода не осознаётся.

Симметрия Румера в настоящее время хорошо известна. Она обсуждается, в частности, в работах Щербака и Макукова. Сделаем новый шаг в анализе этих симметрий. Заметим, что на рис.~1 показан лишь один из возможных способов представления таблицы генетического кода. Другие естественные способы представления можно получить, изменив порядок оснований (T-C-A-G) на сторонах матрицы на другой порядок, \textit{т.\,е.} подвергнув этот ряд перестановке.

Как нетрудно понять, при перестановке оснований на сторонах матрицы столбцы таблицы каким-то образом поменяются местами, и внутри столбцов третье основание кодонов также будет стоять в другом порядке, но полный набор кодирующих продуктов в столбцах останется неизменным. Поэтому, если некоторый столбец принадлежал октету \textbf{I}, он по-прежнему будет принадлежать октету \textbf{I} после любой перестановки оснований на сторонах таблицы, и столбцы октета \textbf{II} также не изменят своей природы. Следовательно, октеты Румера инвариантны относительно преобразований, задаваемых перестановками в строке из четырёх символов: при таких преобразованиях октеты \textbf{I} и \textbf{II} переходят в себя. Любые две последовательные перестановки дают некоторую новую перестановку, существует тождественная перестановка, оставляющая строку оснований неизменной, и каждая перестановка имеет обратную перестановку, возвращающую строку оснований в исходное состояние. Всё это означает, что перестановки образуют группу преобразований,\footnote{Группа~--- это математическая структура, обозначаемая $G$, представляющая собой множество, над элементами которого определена бинарная операция, обычно называемая умножением, и обладающая следующими свойствами: 1) существует единичный элемент $e$ такой, что для любого элемента $\boldsymbol{a}$ из $G$\ $\boldsymbol{a} \cdot \boldsymbol{e} = \boldsymbol{e} \cdot \boldsymbol{a} = \boldsymbol{a}$; 2) для любого элемента $\boldsymbol{a}$ из $G$ существует обратный элемент $\boldsymbol{a}^{-1}$ такой, что $\boldsymbol{a} \cdot \boldsymbol{a}^{-1} = \boldsymbol{a}^{-1} \cdot \boldsymbol{a} = \boldsymbol{e}$; 3) умножение ассоциативно: $\boldsymbol{a} \cdot (\boldsymbol{b} \cdot \boldsymbol{c}) = (\boldsymbol{a} \cdot \boldsymbol{b}) \cdot \boldsymbol{c}$. Основные приложения теории групп связаны с тем, что разнообразные преобразования, такие как вращения фигуры, переводящие фигуру в себя (называемые симметрией этой фигуры), перестановки символов в строке и~\textit{т.\,д.}\ образуют группы. Групповое умножение~--- это последовательное применение двух преобразований, единица группы~--- тождественное преобразование, ничего не меняющее, а обратное преобразование для любого преобразования~--- это то, которое возвращает объект в исходную форму, \textit{напр.}, для вращения фигуры~--- вращение на тот же угол в противоположном направлении и~\textit{т.\,д.}} и группа перестановок четырёх элементов является группой симметрий октетов Румера в том смысле, что октеты Румера не меняются при преобразованиях этой группы.

На рис.~1 видно, что структура областей каллиграммы, соответствующих октетам \textbf{I} и \textbf{II}, довольно сложна. Возникает вопрос: можно ли упростить её путём перестановки порядка оснований на сторонах таблицы так, чтобы области, соответствующие октетам \textbf{I} и \textbf{II}, приняли более простую форму? В частности, на рис.~1 каждая из областей \textbf{I} и \textbf{II} несвязна (более того, каждая из них состоит из трёх отдельных частей). Не существует ли таких расположений оснований на сторонах таблицы генетического кода, при которых каждая из областей, соответствующих октетам \textbf{I} и \textbf{II}, была бы связной, \textit{т.\,е.}\ не имела разрывов?

Для ответа на эти вопросы необходимо проанализировать каллиграммы, соответствующие всем возможным расположениям оснований на сторонах таблицы генетического кода. Всего таких вариантов $4! = 1 \cdot 2 \cdot 3 \cdot 4 = 24$. Рассматривать все 24 каллиграммы не обязательно, поскольку половину из них можно получить инверсией каллиграммы относительно центра квадрата (матрицы), что эквивалентно повороту матрицы на $180^\circ$. Такая операция не меняет фигуру по существу.

\begin{figure}[htbp]
\centering
\includegraphics[width=0.9\textwidth]{fig4_calligrams_top.png}\\[1em]
\includegraphics[width=0.9\textwidth]{fig4_calligrams_bottom.png}

\caption{Все 12 нетривиальных каллиграмм, соответствующих всем возможным представлениям таблицы генетического кода. Ещё 12 каллиграмм получаются поворотом каллиграмм на рисунке на $180^\circ$ (или инверсией относительно центра), что соответствует зеркальному отражению последовательности оснований каждой каллиграммы}
\label{fig:calligrams}
\end{figure}

На рис.~4 показаны все 12 нетривиальных каллиграмм, из которых остальные 12 можно получить поворотом на $180^\circ$. Каллиграмма, соответствующая исходному представлению таблицы на рис.~1, находится в левом верхнем углу рис.~4. Видно, что эта каллиграмма не слишком проста по сравнению с другими. Лишь две каллиграммы особенно просты: CTGA и ATGC. Их простота выражается в том, что каждый из октетов \textbf{I} и \textbf{II} в этих каллиграммах представлен связной областью, то есть непрерывной областью без разрывов.

Оказывается, что все связные каллиграммы (их четыре: две, показанные на рис.~4, и ещё две, получаемые из них поворотом на $180^\circ$) связаны перестановкой оснований, соответствующей преобразованию Румера, и ещё двумя преобразованиями, непосредственно связанными с преобразованием Румера. Наряду с преобразованием Румера $R = (\text{T} \leftrightarrow \text{G}, \text{C} \leftrightarrow \text{A})$ рассмотрим ещё два преобразования, которые в некотором смысле являются <<половинами>> преобразования Румера: $R_1 = (\text{T} \leftrightarrow \text{G})$, $R_2 = (\text{C} \leftrightarrow \text{A})$. В теоретико-групповом смысле $R = R_1 \cdot R_2$, $R_1 = R \cdot R_2$, $R_2 = R \cdot R_1$, и легко проверить, что все три преобразования вместе с тождественным образуют группу (являющуюся подгруппой группы всех перестановок), которая оказывается группой симметрий связных каллиграмм: свойство связности является инвариантом этой группы преобразований. На рис.~5 показано, как все связные каллиграммы переходят друг в друга при преобразованиях $R$, $R_1$, $R_2$ (в математике фигуры такого рода называются коммутативными диаграммами). Мало того что все связные каллиграммы оказываются связанными друг с другом посредством простой группы преобразований, непосредственно связанной с преобразованием Румера, но и сама картинка (см. рис.~5) обладает удивительной симметрией. Легко видеть, что она симметрична относительно инверсии или относительно поворота на $180^\circ$. Вычисления показывают, что вероятность такой дополнительной симметрии при наличии симметрии Румера равна 1/4, то есть полная вероятность получения симметрии Румера вместе с симметрией связных каллиграмм при наличии октетов Румера составляет $1/256 \times 1/4 = 1/1024$.

\begin{figure}[htbp]
\centering
\includegraphics[width=0.5\textwidth]{fig5_commutative.png}
\caption{Действие группы преобразований связных каллиграмм (см. текст)}
\label{fig:commutative}
\end{figure}

Пока мы не могли увидеть никакого скрытого формального смысла в симметриях Румера, всё это можно было интерпретировать как занимательную случайность. Однако эти удивительные симметрии могут служить своего рода <<сигналом привлечения внимания>>: они наводят на мысль о том, что октеты \textbf{I} и \textbf{II} следует рассмотреть более внимательно. Заметим, что на данном этапе даже не столь важно, насколько <<странны>> симметрии Румера в смысле малых вероятностей, благоприятствующих их появлению; важно то, что наличие этих симметрий неизбежно привлекает внимание к октетам Румера как таковым.


\section{Второй информационный уровень. Сигнатуры полных масс кодирующих продуктов}

Принимая во внимание <<подсказку>>, полученную со стороны симметрии Румера, углубим анализ, привлекая новые типы данных. С этой точки зрения вся информация, связанная с симметриями Румера и с дополнительными симметриями связных каллиграмм (см. раздел~3), представляет первый уровень информационных сигнатур генетического кода. Наше продвижение вглубь оправдано вышеупомянутым <<сигналом привлечения внимания>>, но мы по-прежнему будем избегать сложных конструкций и рассматривать только сигнатуры максимальной <<простоты и естественности>>.

Для этого обратим внимание на полные массы кодирующих продуктов (см. таблицу~2). Чтобы подсчитать массы любой группы кодирующих продуктов, мы хотели бы приписать каждому кодирующему продукту некоторый вес, но здесь немедленно возникает нетривиальная проблема. Сигнал \textbf{стоп} тоже является кодирующим продуктом, но он не является материальным объектом, это чисто логическое понятие, ему нельзя естественным образом приписать никакой вес, даже нулевой. Некоторые вещи действительно имеют нулевой вес, например, фотон, но логическая абстракция \textbf{стоп} не имеет веса, что просто не является одним из свойств понятия \textbf{стоп}. Поэтому мы должны искусственно присвоить некоторый вес сигналу \textbf{стоп}. Здесь возникает фундаментальная и нетривиальная неоднозначность (Щербак и Макуков не обратили на неё внимания). Подчеркнём это обстоятельство, поскольку нам придётся воспользоваться им впоследствии. Представляется, что простейшим и наиболее естественным решением является присвоение сигналу \textbf{стоп} веса 0, что мы и сделаем пока. Здесь важно осознать, что мы совершили принципиально важный логический скачок. Суммы масс кодирующих продуктов, с которыми мы будем иметь дело ниже, теперь не являются реальными физическими величинами. Если вычисленная сумма масс включает <<вес>> сигнала \textbf{стоп}, то эта сумма является чисто логическим понятием, и нет смысла даже спрашивать, чем биохимия может объяснить такую сумму. На несколько иную тему аналогичные пояснения даны в работе Щербака и Макукова (2013) (см. раздел~5).

С нашей точки зрения, столь сильный логический шаг был бы совершенно недопустим, если бы анализ начинался сразу с полных масс кодирующих продуктов, без <<сигнала привлечения внимания>> со стороны симметрий Румера. Именно наличие такого <<сигнала>> оправдывает столь нетривиальный логический шаг. Наличие этого особого логического шага позволяет говорить о том, что мы перешли на новый уровень сложности анализа, названный вторым информационным уровнем в заголовке данного раздела. Другим признаком нового информационного уровня является привлечение нового ресурса~--- полных масс кодирующих продуктов.

Теперь заметим, что все кодируемые аминокислоты состоят из так называемой постоянной части и боковой цепи (см. рис.~3). Постоянная часть всех аминокислот, за единственным исключением, одинакова и имеет вес 74, но радикалы все различны. Это единственное исключение~--- аминокислота пролин \textbf{P} (вес постоянной части равен 73), но особая роль числа 74 во всей нашей задаче вполне очевидна, поскольку 19 из 20 аминокислот имеют именно это значение веса постоянной части.

Теперь проделаем очень простую процедуру. Подсчитаем полный вес продуктов, кодируемых в каждом октете по отдельности. При этом мы будем считать каждый продукт ровно столько раз, сколько он встречается в таблице (\textit{т.\,е.} считаем каждую ячейку отдельно). Это представляется простейшим и наиболее естественным способом подсчёта. Например, для столбца CT первого октета продукт \textbf{L} будет посчитан четыре раза, для столбца TG второго октета продукт \textbf{C} будет посчитан дважды, сигнал \textbf{стоп}~--- один раз, а \textbf{W}~--- один раз, и так далее. Это даёт общее число для веса 3700 для первого октета и 4218 для второго октета. По самой своей природе эти числа являются \textit{a priori} случайными целыми числами, поскольку они получены суммированием нескольких целых чисел~--- масс различных аминокислот, не связанных друг с другом. Массы аминокислот не подчиняются никакой простой закономерности (см. таблицу~2).

Оказывается, что каждое из этих двух чисел 3700 и 4218 делится на <<магическое число>> 74, которое уже \textit{a priori} зафиксировано контекстом задачи. Нет никаких оснований ожидать, что полные массы кодирующих продуктов октетов, построенных выше, будут целыми кратными веса постоянной части, заданного \textit{a priori} контекстом задачи, который также является, в определённом смысле, случайным числом. Учитывая случайный характер чисел 3700 и 4218, мы можем оценить вероятность этого странного совпадения. Вероятность того, что случайное целое число делится на 74, равна $1/74$, а вероятность того, что два независимых случайных целых числа оба делятся на 74, составляет всего $1/74 \times 1/74 = 1/5476$. Всё это выглядит, мягко говоря, странно, хотя случайность, разумеется, нельзя исключить. Эта сигнатура не была замечена в работе Щербака и Макукова, поскольку они не работали с полными массами кодирующих продуктов. Их подход с самого начала был более изощрённым (см. раздел~5).

Но это ещё не всё. Вес второго октета оказывается имеющим специальный вид $4218 = 4 \times 999 + 222$. Числа такого вида мы будем называть числами Щербака--Макукова (см. раздел~5). Они играют ведущую роль в информационных сигнатурах последующих уровней.

Но и это ещё не всё. Первый октет Румера проще второго и представляется в некотором смысле более фундаментальным. Если расположить кодирующие продукты первого октета в порядке возрастания веса (что может быть проще?), то первые основания соответствующих кодонов образуют последовательность GGTC|GACC, которая зеркально-симметрична относительно своего центра (отмеченного $|$) в смысле комплементарности;\footnote{Напомним, что при образовании двойной спирали ДНК структура двух соседних цепей связана отношением комплементарности оснований: $\text{G} \leftrightarrow \text{C}$, $\text{T} \leftrightarrow \text{A}$. То есть напротив \textbf{G} всегда стоит \textbf{C}, напротив \textbf{T} всегда стоит \textbf{A}.} нуклеотид C комплементарен нуклеотиду G, нуклеотид T комплементарен нуклеотиду A и так далее. Вероятность, соответствующую этой <<странности>>, не так просто оценить. Здесь возникает частный случай так называемой <<проблемы поиска по углам>> (подробнее см. раздел~9.2.). Применительно к данной ситуации проблема выглядит следующим образом. Несложно подсчитать вероятность появления зеркально-комплементарной симметрии в предположении полной случайности последовательности оснований, она равна 1/256. Но следует учесть, что помимо существующей симметрии есть и другие аналогичные симметрии, которые удивили бы нас не меньше. Например, если бы имело место не зеркальное, а точное повторение первых четырёх букв по комплементарности. Это уже две возможности. Кроме того, можно включить в анализ другие соотношения между основаниями, помимо комплементарности, среди них по крайней мере пары преобразования Румера, полное тождество оснований и транспозиции внутри пуриновой и пиримидиновой пар. Каждое из этих соотношений также включается дважды: для зеркальной симметрии и для точного повторения. В итоге мы получили уже восемь аналогичных симметрий, поэтому полученную выше вероятность следует умножить на 8, получаем $8/256 = 1/32$.

Интересно, что если рассматривать только то, является ли основание в строке GGTCGACC пурином (R) или пиримидином (Y), то, помимо уже существующей комплементарной зеркальной симметрии, мы получаем структуру с симметрией сдвига: RRYYRRYY. Обе эти симметрии были отмечены в статье Щербака и Макукова (2013).

Таким образом, мы спустились на один информационный уровень глубже в сложности анализа и, как и на первом информационном уровне, не остались с пустыми руками. Подчеркнём ещё раз, что для получения числовых и симметрических закономерностей второго уровня сложности нам не пришлось проделывать никаких сложных манипуляций с цифрами, кроме искусственного присвоения нулевого веса сигналу \textbf{стоп} (что на первый взгляд кажется очень простым). Все результаты <<лежат на поверхности>> и получаются естественным образом. Фактически информационные закономерности второго уровня настолько просты, что можно сказать, что оба информационных уровня играют роль <<сигналов привлечения внимания>> в определённом смысле. Однако со вторым информационным уровнем мы получили значительное усиление <<сигнала привлечения внимания>>, что позволяет начать поиск более тонких и сложных информационных закономерностей генетического кода. Это даёт моральное оправдание для более сложных манипуляций с обсуждаемыми данными.


\section{Третий информационный уровень. Массы боковых цепей}

Основной массив информационных сигнатур статьи (Щербак и Макуков 2013) был получен не с использованием полных масс аминокислот, как мы получили массы октетов \textbf{I} и \textbf{II} (3700 и 4218) на втором информационном уровне, а с использованием масс боковых цепей аминокислот. Здесь, однако, обнаруживается одно удивительное обстоятельство. Значительная часть информационных закономерностей появляется лишь в том случае, если искусственно скорректировать структуру молекулы пролина, имеющего нестандартный вес постоянной части: 73 вместо 74. Необходимо искусственно (виртуально!) перенести один водород из боковой цепи пролина в постоянную часть, после чего вес боковой цепи становится 41 вместо исходных 42, а вес постоянной части становится стандартным, то есть 74. Щербак и Макуков называют эту мысленную операцию <<ключом активации>>, дающим доступ к основному массиву сигнатур генетического кода. Как справедливо отмечено в работе Щербака и Макукова (2013), все сигнатуры приобретают виртуальный характер: различные суммы масс боковых цепей аминокислот с учётом поправки на пролин не являются чем-то реальным, поэтому вопрос о том, почему эти суммы именно такие, а не другие, теряет свой физический смысл.

Проблема, аналогичная уже обсуждавшейся выше, вновь возникает с кодирующим продуктом \textbf{стоп}. Если мы теперь рассматриваем не полные массы аминокислот, а только массы боковых цепей, какой вес <<боковой цепи>> следует приписать сигналу \textbf{стоп}? Щербак и Макуков искусственно присваивают нулевой вес сигналу \textbf{стоп}, и именно этот и только этот выбор приводит к большому набору информационных закономерностей в статье (\textit{Там же}.). По существу, используется не один <<ключ активации>>, как полагают Щербак и Макуков, а два. Переход к использованию масс боковых цепей аминокислот вместе с двумя искусственными корректировками знаменует переход к третьему информационному уровню сложности.

Искусственное присвоение нулевого веса <<боковой цепи>> сигнала \textbf{стоп} имеет некоторые интересные следствия, которые не были рассмотрены Щербаком и Макуковым. После стандартизации веса постоянной части все массы боковых цепей аминокислот вычисляются как полный вес молекулы минус 74. Для единообразного поведения всех кодирующих продуктов следовало бы предположить, что то же самое справедливо и для сигнала \textbf{стоп}. Но для того чтобы получить ноль для веса <<боковой цепи>> сигнала \textbf{стоп}, мы должны предположить, что <<полный вес>> сигнала \textbf{стоп} равен не нулю, как мы предполагали на втором информационном уровне (см. раздел~4), а 74. Только в этом случае мы получаем правильное значение для <<боковой цепи>> сигнала \textbf{стоп}: $74 - 74 = 0$. Возможность переопределения веса сигнала \textbf{стоп} согласуется со свободой выбора полного веса продукта \textbf{стоп}, о которой мы писали выше (см. раздел~4). Хотя выбор нулевого значения полного веса сигнала \textbf{стоп} представляется простейшим, он в определённом смысле не является наиболее логичным (или правильным).

С этим новым пониманием природы <<полного веса>> сигнала \textbf{стоп} мы можем вернуться к вычислению полных масс октетов Румера (см. раздел~4). Для первого октета с его полным весом 3700 ничего не меняется, поскольку сигнал \textbf{стоп} не входит в подсчёт его веса, но вес второго октета меняется с 4218 на 4440. Именно здесь происходит маленькое вероятностное чудо. В то время как числа 3700 и 4218 имели наибольший общий делитель 74, что уже было неожиданно большим числом, числа 3700 и 4440 имеют наибольший общий делитель 740! Оказывается, что присвоение полного веса 74 сигналу \textbf{стоп} не только соответствует <<ключу активации>> сигнатур, открытых Щербаком и Макуковым, но и существенно усиливает сигнатуру второго информационного уровня~--- одного из двух <<уровней привлечения внимания>>. Но, оказывается, это ещё не всё.

Разложим новый наибольший общий делитель 740 на множители следующим образом:
\begin{equation}
740 = 2 \times 10 \times 37.
\end{equation}
Теперь заметим, что в нашей задаче с самого начала выделены два числа~--- вес постоянной части аминокислот 74 и число аминокислот, встречающихся в генетическом коде, 20. Запишем эти числа следующим образом:
\begin{equation}
20 = 2 \times 10; \quad 74 = 2 \times 37.
\end{equation}
Во всех числах в правых частях уравнений (1, 2) присутствует общий множитель 2, на который мы сокращаем эти числа, после чего получаем следующий ряд чисел: $10 \times 37$, 10, 37. Пара чисел (10, 37) появляется двумя независимыми способами сразу: первый раз как сомножители числа 370, второй раз по отдельности. Это может указывать на какую-то особую роль этой пары чисел.\footnote{Заметим, что $3700 = 37 \times 10 \times 10$, что тоже намекает на особую роль чисел 37 и 10, но 10 встречается в этом разложении дважды, поэтому мы не рассматриваем эту <<сигнатуру>> как прямое указание на пару (10, 37).}

Можно было бы сказать, что всё это не более чем упражнение в нумерологии, если бы эта пара чисел не играла абсолютно исключительную роль в генерации всего основного набора сигнатур генетического кода, найденных Щербаком и Макуковым. Более того, рассмотрение особых свойств пары (10, 37) является необходимым прологом к обзору этих сигнатур.

Начать следует с того, что пара чисел (10, 37) обладает некоторыми замечательными математическими свойствами,\footnote{Здесь наше изложение следует работе (Щербак и Макуков 2013) с минимальными вариациями.} которые были открыты Пачоли в 1508 году и известны давно.

Запишем в таблицу результаты умножения числа 37 на натуральные числа от 1 до 54 (см. таблицу~3). Если посмотреть на трёхзначные результаты умножения, можно увидеть, что среди них встречаются все числа вида $n \times 111$, где $n$ изменяется от 1 до 9 (выделены жирным шрифтом). Более того, для всех этих чисел сумма цифр результата равна множителю числа 37, с которым этот результат получен: $3 + 3 + 3 = 9$, $6 + 6 + 6 = 18$ и так далее. Для всех остальных трёхзначных результатов умножения результат каждого столбца получается из результата верхней строки циклической перестановкой цифр, например $2 \times 37 = 074$, $11 \times 37 = 407$, $20 \times 37 = 740$. Далее, для множителей от 28 до 54 закономерность повторяется полностью, но ко всем результатам умножения прибавляется слагаемое $1 \times 999$. В следующем цикле появляется слагаемое $2 \times 999$ и так далее. Причём во всех результатах вида $m \times 999 + n \times 111$ множитель, с помощью которого получено число, по-прежнему равен сумме цифр результата в определённом смысле, а именно, он равен $m \times (9 + 9 + 9) + n \times (1 + 1 + 1)$. Числа вида $m \times 999 + n \times 111$ играют ключевую роль в информационных сигнатурах генетического кода. Это открытие Щербака и Макукова. Числа такого вида мы будем называть числами Щербака--Макукова, или сокращённо SM-числами. Таким образом, простейшие SM-числа имеют вид $n \times 111$ при $n = 1, \ldots, 9$, за ними следуют числа вида $1 \times 999 + n \times 111$ и так далее.

\begin{table}[htbp]
\centering
\caption{Свойства симметрии таблицы умножения числа 37 в позиционной системе счисления с основанием 10}
\label{tab:mult37}
\footnotesize
\begin{tabular}{|c|c|c|c|c|c|c|c|c|}
\hline
1 & 2 & 3 & 4 & 5 & 6 & 7 & 8 & 9 \\
037 & 074 & \textbf{111} & 148 & 185 & \textbf{222} & 259 & 296 & \textbf{333} \\
\hline
10 & 11 & 12 & 13 & 14 & 15 & 16 & 17 & 18 \\
370 & 407 & \textbf{444} & 481 & 518 & \textbf{555} & 592 & 629 & \textbf{666} \\
\hline
19 & 20 & 21 & 22 & 23 & 24 & 25 & 26 & 27 \\
703 & 740 & \textbf{777} & 814 & 851 & \textbf{888} & 925 & 962 & \textbf{999} \\
\hline
28 & 29 & 30 & 31 & 32 & 33 & 34 & 35 & 36 \\
999+037 & 999+074 & 999+\textbf{111} & 999+148 & 999+185 & 999+\textbf{222} & 999+259 & 999+296 & 999+\textbf{333} \\
\hline
37 & 38 & 39 & 40 & 41 & 42 & 43 & 44 & 45 \\
999+370 & 999+407 & 999+\textbf{444} & 999+481 & 999+518 & 999+\textbf{555} & 999+592 & 999+629 & 999+\textbf{666} \\
\hline
46 & 47 & 48 & 49 & 50 & 51 & 52 & 53 & 54 \\
999+703 & 999+740 & 999+\textbf{777} & 999+814 & 999+851 & 999+\textbf{888} & 999+925 & 999+962 & 999+\textbf{999} \\
\hline
\end{tabular}
\end{table}

Заметим, что по крайней мере одно SM-число можно найти уже на втором информационном уровне, без использования масс боковых цепей аминокислот. Как бы мы ни вычисляли вес второго октета Румера~--- в предположении нулевого веса сигнала \textbf{стоп} или в предположении веса, равного 74,~--- результат оказывается числом Щербака--Макукова:
\begin{align*}
M_{stop} = 0 &\Rightarrow 4218 = 4 \times 999 + 222 \\
M_{stop} = 74 &\Rightarrow 4440 = 4 \times 999 + 444.
\end{align*}
Удивительная элегантность последнего варианта может быть дополнительным указанием на предпочтительность выбора <<полного веса>> 74 для сигнала \textbf{стоп} (эти сигнатуры не были замечены Щербаком и Макуковым).

На первый взгляд таблица~3 представляет особые свойства числа 37, но на самом деле эта любопытная арифметика получается только в позиционной системе счисления с основанием 10, то есть это свойства пары чисел (10, 37). Пара чисел (10, 37) не уникальна в этом отношении. Любая пара целых чисел $(q, 111_q / 3)$, где $q = 4, 7, 10, 13, \ldots$ и $111_q$ означает 111 в позиционной системе счисления с основанием $q$ (а именно: $111_q = 1 \times q^2 + 1 \times q + 1$), обладает аналогичными свойствами (помимо приведённой выше таблицы умножения, у них есть и некоторые другие любопытные свойства, которых мы здесь не касались [см. Щербак и Макуков 2013]). Это означает, что, хотя свойства пары (10, 37) интересны, не следует переоценивать степень необычности этих свойств.

Начиная с этого момента мы будем использовать другой способ подсчёта масс по группам кодонов, нежели тот, который мы использовали на втором информационном уровне (см. раздел~4), когда мы считали каждую ячейку таблицы генетического кода отдельно. Проходя по таблице генетического кода различными способами (см. ниже), мы будем считать вес каждого кодирующего продукта (или вес боковой цепи, соответствующей этому продукту) для каждого столбца таблицы, в котором он встречается, только один раз, не учитывая степень вырожденности этого продукта в данном столбце. Это правило подсчёта довольно нетривиально (предложено и использовано Щербаком и Макуковым), и для него нет \textit{априорного} обоснования. Однако суть в том, что приводимые ниже информационные сигнатуры получаются именно этим, и никаким другим способом, так что этот метод подсчёта обоснован \textit{апостериорно}.

Перечислим теперь информационные структуры генетического кода, полученные с использованием масс боковых цепей аминокислот (третий информационный уровень). Большинство из них имеет вид чисел Щербака--Макукова или связано с такими числами. Таких структур много.

\subsection*{Сигнатура 3.1. Боковые цепи первого и второго октета и египетский треугольник}

Начнём с первого октета Румера, который, как отмечалось выше, представляется более фундаментальным, чем второй. Сумма масс боковых цепей аминокислот по правилу подсчёта Щербака--Макукова и с учётом <<ключа активации>> пролина (искусственное переопределение веса боковой цепи и постоянной части для пролина $42 \to 41$, $73 \to 74$) есть SM-число 333 (см. таблицу~4):
\begin{equation}
\text{боковые цепи октета \textbf{I}: } 333 = 37 \times 9.
\end{equation}

\begin{table}[htbp]
\centering
\caption{Суммы масс аминокислот октета I}
\label{tab:octet1}
\begin{tabular}{|l|c|c|c|c|c|c|c|c|c|}
\hline
аминокислота & G & A & S & P & V & T & L & R & $\Sigma$ \\
\hline
боковая цепь & 1 & 15 & 31 & 41 & 43 & 45 & 57 & 100 & \textbf{333} \\
\hline
постоянная часть & 74 & 74 & 74 & 74 & 74 & 74 & 74 & 74 & \textbf{592} \\
\hline
полный вес & 75 & 89 & 105 & 115 & 117 & 119 & 131 & 174 & \textbf{925} \\
\hline
\end{tabular}
\quad
\raisebox{-0.5\height}{\includegraphics[width=2cm]{cal_tab4_octetI.png}}
\end{table}

Если теперь выписать сумму масс постоянных частей аминокислот и полные массы аминокислот октета \textbf{I}, мы получим, соответственно:
\begin{equation}
\text{постоянные части октета \textbf{I}: } 592 = 37 \times 16
\end{equation}
\begin{equation}
\text{полные массы октета \textbf{I}: } 925 = 37 \times 25.
\end{equation}
Все числа в уравнениях (3--5) кратны 37, и если мы сократим их на этот общий множитель, то получим три последовательных квадрата целых чисел $3^2$, $4^2$, $5^2$, которые, как трудно не заметить, вместе образуют египетский треугольник: $3^2 + 4^2 = 5^2$. Можно также отметить, что это представление теоремы Пифагора появляется не в каком-то случайном месте, а именно в первом октете Румера, который напоминает структурный центр генетического кода (и фактически является его наиболее стабильной частью), причём эта сигнатура не требует сложного правила для выбора группы аминокислот, для которой считается полный вес. Как будет видно ниже, остальные правила выбора групп для подсчёта масс боковых цепей, хотя всегда имеют некоторый \textit{априорный} смысл, более сложны. Следует также отметить, что эта сигнатура свободна от неопределённости, связанной со свободой выбора веса сигнала \textbf{стоп}, хотя она критически зависит от <<ключа активации>>, связанного с переопределением структуры пролина. Использование <<ключа активации>>, как уже отмечалось, оправдано \textit{апостериорно} тем, что оно приводит не к одной сигнатуре, а ко множеству независимых сигнатур одновременно.

Используя правило подсчёта Щербака--Макукова для суммы масс боковых цепей кодирующих продуктов октета \textbf{II}, мы получаем сумму масс, которая также является SM-числом: $1110 = 999 \times 1 + 111$ (см. таблицу~5). Напомним, что мы уже один раз получили SM-число для октета \textbf{II}: это был полный вес октета \textbf{II}, полученный подсчётом каждой ячейки таблицы генетического кода по отдельности. Эти две сигнатуры независимы друг от друга.

\begin{table}[htbp]
\centering
\caption{Суммы масс аминокислот октета II. Порядок представления аминокислот в таблице соответствует движению слева направо в каждой строке генетического кода (см. рис.~1) и сверху вниз в каждом столбце генетического кода}
\label{tab:octet2}
\begin{tabular}{|l|c|c|c|c|c|c|c|c|c|}
\hline
аминокислота & F & L & Y & \textbf{стоп} & C & \textbf{стоп} & W & H & Q $\rightarrow$ \\
\hline
боковая цепь & 91 & 57 & 107 & 0 & 47 & 0 & 130 & 81 & $72 \rightarrow$ \\
\hline
постоянная часть & 74 & 74 & 74 & 0 & 74 & 0 & 74 & 74 & $74 \rightarrow$ \\
\hline
аминокислота & I & M & N & K & S & R & D & E & $\Sigma$ \\
\hline
боковая цепь & 57 & 75 & 58 & 72 & 31 & 100 & 59 & 73 & \textbf{999+111} \\
\hline
постоянная часть & 74 & 74 & 74 & 74 & 74 & 74 & 74 & 74 & \textbf{999+111} \\
\hline
\end{tabular}
\quad
\raisebox{-0.5\height}{\includegraphics[width=2cm]{cal_tab5_octetII.png}}
\end{table}

Вероятно, существует ещё одна сигнатура, связанная с октетом \textbf{II}. Если мы предположим, что полный вес сигнала \textbf{стоп} равен нулю (а значит, и вес постоянной части сигнала \textbf{стоп} равен нулю), то полный вес постоянных частей октета \textbf{II} при подсчёте по правилу Щербака--Макукова даёт SM-число $999 \times 1 + 111$. Другими словами, имеется дополнительное равновесие масс боковых и постоянных частей октета \textbf{II}: $999 \times 1 + 111 = 999 \times 1 + 111$. При более <<разумном>> предположении $M_{stop} = 74$ это равновесие нарушается.


\subsection*{Сигнатура 3.2. Кодоны с двумя одинаковыми основаниями}

Существует 36 кодонов, содержащих два одинаковых основания и ещё одно основание, отличное от этих двух. Эту группу оснований можно разделить на две половины таким образом, что в одной группе будут кодоны, в которых пара одинаковых оснований~--- пиримидины, а в другой~--- пурины. Тогда суммы масс каждой из половин равны (равновесие масс) и равны числу Щербака--Макукова 999 (см. таблицы~6 и 7). Это первый пример выбора групп аминокислот для суммирования масс, который имеет \textit{априорный} смысл, но который довольно сложен по сравнению с сигнатурой~3.1.

\begin{table}[htbp]
\centering
\caption{Масса боковых цепей кодонов, в которых одна пара одинаковых оснований~--- пиримидины, а другое основание отлично от них}
\label{tab:pyrimidine_pairs}
\begin{tabular}{|c|c|c|c|c|c|}
\hline
TTC & TTA & TTG & CCT & CCA & CCG \\
\textbf{F} 91 & \textbf{L} 57 & \textbf{L} 57 & \textbf{P} 41 & \textbf{P} 41 & \textbf{P} 41 \\
\hline
CTT & ATT & GTT & TCC & ACC & GCC \\
\textbf{L} 57 & \textbf{I} 57 & \textbf{V} 43 & \textbf{S} 31 & \textbf{T} 45 & \textbf{A} 15 \\
\hline
TCT & TAT & TGT & CTC & CAC & CGC \\
\textbf{S} 31 & \textbf{Y} 107 & \textbf{C} 47 & \textbf{L} 57 & \textbf{H} 81 & \textbf{R} 100 \\
\hline
\multicolumn{6}{|c|}{$\Sigma = 999$} \\
\hline
\end{tabular}
\quad
\raisebox{-0.5\height}{\includegraphics[width=2cm]{cal_tab6_pyrimidine.png}}
\end{table}

\begin{table}[htbp]
\centering
\caption{Масса боковых цепей кодонов, в которых одна пара одинаковых оснований~--- пурины, а другое основание отлично от них. Различные цвета заливки на диаграмме генетического кода соответствуют различным подгруппам кодонов в таблице, разделённым двойной линией}
\label{tab:purine_pairs}
\begin{tabular}{|c|c|c||c|c|c|}
\hline
 & AAT & & GGT & GGC & GGA \\
\multirow{2}{*}{N 58} & AAC & AAG & \textbf{G} 1 & \textbf{G} 1 & \textbf{G} 1 \\
 & \textbf{N} 58 & \textbf{K} 72 & & & \\
\hline
TAA & CAA & GAA & TGG & CGG & AGG \\
\textbf{стоп} 0 & \textbf{Q} 72 & \textbf{E} 73 & \textbf{W} 130 & \textbf{R} 100 & \textbf{R} 100 \\
\hline
ATA & ACA & AGA & GTG & GCC & GAG \\
\textbf{I} 57 & \textbf{T} 45 & \textbf{R} 100 & \textbf{V} 43 & \textbf{A} 15 & \textbf{E} 73 \\
\hline
\multicolumn{6}{|c|}{$\Sigma = 999$} \\
\hline
\end{tabular}
\quad
\raisebox{-0.5\height}{\includegraphics[width=2cm]{cal_tab7_purine.png}}
\end{table}

Щербак и Макуков показали, что существует вариант разбиения полностью однородных и полностью неоднородных групп пополам, при котором для сумм получаются некоторые SM-числа, но это специальный выбор среди многих в остальном равноправных вариантов, не основанный на каком-либо \textit{априорно} ясном правиле (вроде пуринов с одной стороны, пиримидинов с другой), поэтому этот результат может считаться результатом подгонки. Мы его здесь не приводим.


\subsection*{Сигнатура 3.3. Тройная симметрия кодонов с двумя одинаковыми основаниями, являющимися пуринами}

Группу кодонов, в которой пара одинаковых оснований~--- пурины (с весом 999, см. сигнатуру~3.2.\ и таблицу~7), можно разделить на три равные подгруппы по следующему принципу: первая группа содержит два основания A подряд, вторая группа содержит два основания G подряд, третья группа содержит либо два основания A, разделённых другим основанием, либо два основания G, разделённых другим основанием. Эти три группы разделены двойными линиями в таблице~7. Оказывается, что сумма масс боковых цепей каждой подгруппы равна SM-числу 333. Другими словами, имеется тройное равновесие подгрупп, и число равновесия является SM-числом.


\subsection*{Сигнатура 3.4. Кодоны с двумя одинаковыми основаниями и единственным пурином или единственным пиримидином}

Если в группе кодонов с двумя одинаковыми основаниями и отличным третьим основанием (как в сигнатуре~3.2.) рассмотреть отдельно кодоны с единственным пурином и единственным пиримидином, оказывается, что соответствующие суммы масс боковых цепей являются SM-числами по отдельности: для единственного пурина 888 (см. таблицу~8); для единственного пиримидина $1 \times 999 + 111$ (см. таблицу~9). Независимой сигнатурой здесь является свойство SM-числа каждой суммы в отдельности.

\begin{table}[htbp]
\centering
\caption{Два одинаковых основания и отличный третий единственный пурин}
\label{tab:single_purine}
\begin{tabular}{|c|c|c|c|c|c|}
\hline
TTC & CCT & AAT & AAC & GGT & GGC \\
\textbf{F} 91 & \textbf{P} 41 & \textbf{N} 58 & \textbf{N} 58 & \textbf{G} 1 & \textbf{G} 1 \\
\hline
CTT & TCC & TAA & CAA & TGG & CGG \\
\textbf{L} 57 & \textbf{S} 31 & \textbf{стоп} 0 & \textbf{Q} 72 & \textbf{W} 130 & \textbf{R} 100 \\
\hline
TCT & CTC & ATA & ACA & GTG & GCG \\
\textbf{S} 31 & \textbf{L} 57 & \textbf{I} 57 & \textbf{T} 45 & \textbf{V} 43 & \textbf{A} 15 \\
\hline
\multicolumn{6}{|c|}{$\Sigma = 888$} \\
\hline
\end{tabular}
\quad
\raisebox{-0.5\height}{\includegraphics[width=2cm]{cal_tab8_single_purine.png}}
\end{table}

Таблица~9 построена по простому принципу симметрии. Первая строка содержит кодоны с единственным пиримидином в последней позиции, где первые два основания одинаковы. Следующие две строки получены циклической перестановкой позиции единственного пиримидина по часовой стрелке. Сумма второй строки равна 333, что является независимой SM-сигнатурой. Итого $1 \times 999 + 111 = 333 + 777$.

\begin{table}[htbp]
\centering
\caption{Два одинаковых основания и отличный третий единственный пиримидин. Группа кодонов с суммой масс 333 выделена чёрным цветом на диаграмме генетического кода}
\label{tab:single_pyrimidine}
\begin{tabular}{|c|c|c|c|c|c|c|}
\hline
TTA & TTG & CCA & CCG & AAG & GGA & $\Sigma$ \\
\textbf{L} 57 & \textbf{L} 57 & \textbf{P} 41 & \textbf{P} 41 & \textbf{K} 72 & \textbf{G} 1 & 269 \\
\hline
ATT & GTT & ACC & GCC & GAA & AGG & \\
\textbf{I} 57 & \textbf{V} 43 & \textbf{T} 45 & \textbf{A} 15 & \textbf{E} 73 & \textbf{R} 100 & 333 \\
\hline
TAT & TGT & CAC & CGC & AGA & GAG & \\
\textbf{Y} 107 & \textbf{C} 47 & \textbf{H} 81 & \textbf{R} 100 & \textbf{R} 100 & \textbf{E} 73 & 508 \\
\hline
\multicolumn{7}{|c|}{$\Sigma = 1 \times 999 + 111 = 333 + 777$} \\
\hline
\end{tabular}
\quad
\raisebox{-0.5\height}{\includegraphics[width=2cm]{cal_tab9_single_pyrimidine.png}}
\end{table}


\subsection*{Сигнатура 3.5. Полные массы кодонов с первым пиримидином или первым пурином}

Если разделить все 64 кодона на две группы по тому, является ли первое основание пиримидином или пурином, и подсчитать сумму полных масс кодирующих продуктов (без учёта вырожденности, с весом \textbf{стоп} = 74), то сумма для кодонов с первым пиримидином равна $1 \times 999 + 777$ (см. таблицу~10), а сумма для кодонов с первым пурином равна $1517 = 41 \times 37$ (см. таблицу~11). Последнее не является SM-числом, но кратно 37. Эти две сигнатуры отсутствовали в работе Щербака и Макукова (2013).

\begin{table}[htbp]
\centering
\caption{Полный вес кодирующих продуктов, соответствующих кодонам с первым пиримидином}
\label{tab:first_pyrimidine}
\begin{tabular}{|c|c|c|c|c|c|c|}
\hline
TTY & TTR & TCN & TAY & TAR & TGY & \\
\textbf{F} 165 & \textbf{L} 131 & \textbf{S} 105 & \textbf{Y} 181 & \textbf{стоп} 74 & \textbf{C} 121 & \\
\hline
TGA & TGG & CTN & CCN & CAY & CAR & CGN \\
\textbf{стоп} 74 & \textbf{W} 204 & \textbf{L} 131 & \textbf{P} 115 & \textbf{H} 155 & \textbf{Q} 146 & \textbf{R} 174 \\
\hline
\multicolumn{7}{|c|}{$\Sigma = 1 \times 999 + 777$} \\
\hline
\end{tabular}
\quad
\raisebox{-0.5\height}{\includegraphics[width=2cm]{cal_tab10_first_pyrimidine.png}}
\end{table}

\begin{table}[htbp]
\centering
\caption{Полный вес кодирующих продуктов, соответствующих кодонам с первым пурином}
\label{tab:first_purine}
\begin{tabular}{|c|c|c|c|c|c|c|}
\hline
ATH & ATG & ACN & AAY & AAR & AGY & \\
\textbf{I} 131 & \textbf{M} 149 & \textbf{T} 119 & \textbf{N} 132 & \textbf{K} 146 & \textbf{S} 105 & \\
\hline
AGR & GTN & GCN & GAY & GAR & GGN & \\
\textbf{R} 174 & \textbf{V} 117 & \textbf{A} 89 & \textbf{D} 133 & \textbf{E} 147 & \textbf{G} 75 & \\
\hline
\multicolumn{7}{|c|}{$\Sigma = 1517 = 41 \times 37$} \\
\hline
\end{tabular}
\quad
\raisebox{-0.5\height}{\includegraphics[width=2cm]{cal_tab11_first_purine.png}}
\end{table}


\subsection*{Сигнатура 3.6. Разложение кодонов}

Каждый кодон можно мысленно разложить на составляющие его основания, приписав каждому основанию вес продукта, кодируемого исходным кодоном. Сумма масс боковых цепей на всех отдельных основаниях с учётом вырожденности есть SM-число $10 \times 999 + 222$.

Эта сигнатура не является независимой~--- она может быть выведена из октетов ($3700 + 4440 = 8140$, минус $74 \times 64$ постоянных частей $= 3404$, также делится на 37; умножаем на 3 для отдельных оснований: $37 \times 3 = 111$, так что обязательно SM-число: $3404 \times 3 = 10 \times 999 + 222$).

Однако сумма для основания T отдельно равна SM-числу $2 \times 999 + 666$, что является независимой сигнатурой. Сумма для C, G, A вместе равна $7 \times 999 + 555$, что следует из того, что разность SM-чисел есть SM-число.


\subsection*{Сигнатура 3.7. Пары Румера M = (T, G) и K = (C, A)}

Если разделить столбцы таблицы генетического кода по тому, принадлежит ли первое основание паре Румера $\text{M} = (\text{T}, \text{G})$ или $\text{K} = (\text{C}, \text{A})$, мы получим следующие массы боковых цепей.

\begin{table}[htbp]
\centering
\caption{Сумма масс боковых цепей кодонов с первым основанием, принадлежащим паре M = (T, G)}
\label{tab:halfM_first}
\begin{tabular}{|c|c|c|c|c|c|c|c|c|c|}
\hline
аминокислота & F & L & Y & \textbf{стоп} & C & \textbf{стоп} & W & S & $\rightarrow$ \\
\hline
боковая цепь & 91 & 57 & 107 & 0 & 47 & 0 & 130 & 31 & $\rightarrow$ \\
\hline
аминокислота & V & A & D & E & G & & & & $\Sigma$ \\
\hline
боковая цепь & 43 & 15 & 59 & 73 & 1 & & & & \textbf{654} \\
\hline
\end{tabular}
\quad
\raisebox{-0.5\height}{\includegraphics[width=2cm]{cal_tab12_M_first.png}}
\end{table}

\begin{table}[htbp]
\centering
\caption{Сумма масс боковых цепей кодонов с первым основанием, принадлежащим паре K = (C, A)}
\label{tab:halfK_first}
\begin{tabular}{|c|c|c|c|c|c|c|c|c|c|}
\hline
аминокислота & L & P & H & Q & R & I & M & N & $\rightarrow$ \\
\hline
боковая цепь & 57 & 41 & 81 & 72 & 100 & 57 & 75 & 58 & $\rightarrow$ \\
\hline
аминокислота & K & S & R & T & & & & & $\Sigma$ \\
\hline
боковая цепь & 72 & 31 & 100 & 45 & & & & & \textbf{789} \\
\hline
\end{tabular}
\quad
\raisebox{-0.5\height}{\includegraphics[width=2cm]{cal_tab13_K_first.png}}
\end{table}

Отдельные суммы 654 и 789 не делятся на 37 и по отдельности не интересны. Однако их сумма $789 + 654 = 1 \times 999 + 444$, что является SM-числом, но это не является независимой сигнатурой, поскольку следует из сигнатуры~3.1.

Если вместо первого основания использовать для группировки среднее (второе) основание кодона, мы получим: группа~M $= 789$ (таблица~14), группа~K $= 654$ (таблица~15)~--- те же значения с точностью до перестановки, и алгебраически независимые.

\begin{table}[htbp]
\centering
\caption{Сумма масс боковых цепей кодонов со вторым основанием, принадлежащим паре M}
\label{tab:groupM_second}
\begin{tabular}{|c|c|c|c|c|c|c|c|c|c|}
\hline
аминокислота & F & L & L & I & M & V & C & $\rightarrow$ & \\
\hline
боковая цепь & 91 & 57 & 57 & 57 & 75 & 43 & 47 & $\rightarrow$ & \\
\hline
аминокислота & \textbf{стоп} & W & R & S & R & G & & $\Sigma$ & \\
\hline
боковая цепь & 0 & 130 & 100 & 31 & 100 & 1 & & \textbf{789} & \\
\hline
\end{tabular}
\quad
\raisebox{-0.5\height}{\includegraphics[width=2cm]{cal_tab14_M_second.png}}
\end{table}

\begin{table}[htbp]
\centering
\caption{Сумма масс боковых цепей кодонов со вторым основанием, принадлежащим паре K}
\label{tab:groupK_second}
\begin{tabular}{|c|c|c|c|c|c|c|c|c|c|}
\hline
аминокислота & S & P & T & A & Y & \textbf{стоп} & H & Q & $\rightarrow$ \\
\hline
боковая цепь & 31 & 41 & 45 & 15 & 107 & 0 & 81 & 72 & $\rightarrow$ \\
\hline
аминокислота & N & K & D & E & & & & & $\Sigma$ \\
\hline
боковая цепь & 58 & 72 & 59 & 73 & & & & & \textbf{654} \\
\hline
\end{tabular}
\quad
\raisebox{-0.5\height}{\includegraphics[width=2cm]{cal_tab15_K_second.png}}
\end{table}

Далее, если рассмотреть оба первых кодона в M и оба первых кодона в K, мы получим баланс 369 = 369 (таблица~16), который алгебраически независим.

\begin{table}[htbp]
\centering
\caption{Массы групп кодирующих продуктов, у которых оба первых основания принадлежат группам M или K}
\label{tab:balance_MK}
\begin{tabular}{|c|c|c|c||c|c|c|}
\hline
\multicolumn{4}{|c||}{Оба первых основания в M} & \multicolumn{3}{c|}{Оба первых основания в K} \\
\hline
TTY & TTR & TGY & TGA & CCN & CAY & CAR \\
\textbf{F} 91 & \textbf{L} 57 & \textbf{C} 47 & \textbf{стоп} 0 & \textbf{P} 41 & \textbf{H} 81 & \textbf{Q} 72 \\
\hline
TGG & GTN & GGN & & ACN & AAY & AAR \\
\textbf{W} 130 & \textbf{V} 43 & \textbf{G} 1 & & \textbf{T} 45 & \textbf{N} 58 & \textbf{K} 72 \\
\hline
\multicolumn{4}{|c||}{$\Sigma = 369$} & \multicolumn{3}{c|}{$\Sigma = 369$} \\
\hline
\end{tabular}
\quad
\raisebox{-0.5\height}{\includegraphics[width=2cm]{cal_tab16_MK_both.png}}
\end{table}


\section{Третий информационный уровень (продолжение). TGA-переключатель}

Стандартный генетический код содержит элемент, нарушающий некоторые из обсуждавшихся выше симметрий: кодон TGA кодирует \textbf{стоп}, тогда как <<симметричное>> ожидание состояло бы в том, что он кодирует цистеин (\textbf{C}). Рассмотрим симметризованный генетический код:
\begin{equation}
\text{TGA} : \textbf{стоп} \to \text{TGA} : \textbf{C}.
\end{equation}
Эта симметризация существует в природе как альтернативный код эуплотидиума.

Симметризация разрушает сигнатуры второго информационного уровня, но создаёт новые на третьем уровне. Мы предлагаем рассматривать это как двухпозиционный переключатель:
\begin{equation}
\text{TGA} : \textbf{стоп} \leftrightarrow \text{TGA} : \textbf{C}.
\end{equation}

\begin{table}[htbp]
\centering
\caption{Идеограмма второго октета симметризованного генетического кода (пояснение в тексте)}
\label{tab:ideogram_symm}
\footnotesize
\begin{tabular}{|c|c|c|c|c|c|c|c|c|c|c|c|c|c|c|c|c|}
\hline
группа & \multicolumn{2}{c|}{\textbf{III}} & \multicolumn{12}{c|}{\textbf{II}} & \multicolumn{2}{c|}{\textbf{I}} \\
\hline
продукт & \textbf{C} & \textbf{I} & \textbf{стоп} & \textbf{S} & \textbf{L} & \textbf{N} & \textbf{D} & \textbf{Q} & \textbf{K} & \textbf{E} & \textbf{H} & \textbf{F} & \textbf{R} & \textbf{Y} & \textbf{M} & \textbf{W} \\
$M_{\text{side}}$ & 47 & 57 & 0 & 31 & 57 & 58 & 59 & 72 & 72 & 73 & 81 & 91 & 100 & 107 & 75 & 130 \\
\hline
основание 1 & T & A & T & A & T & A & G & C & A & G & C & T & A & T & A & T \\
основание 2 & G & T & A & G & T & A & A & A & A & A & A & T & G & A & T & G \\
основание 3 & H & H & R & Y & R & Y & Y & R & R & R & Y & Y & R & Y & G & G \\
\hline
\end{tabular}
\end{table}

\paragraph{Сигнатура 3.1B: Симметрии строки основания 1.} В симметризованном коде первые и последние пять оснований строки основания~1 симметричны (паттерн TATAT), а центральный сегмент AGC$|$AGC зеркально симметричен.

\paragraph{Сигнатура 3.2B: Строка основания 2.} Строка основания~2 симметризованного кода точно зеркально симметрична.

\paragraph{Сигнатура 3.3B: Симметрия между группами вырожденности.}

\begin{table}[htbp]
\centering
\caption{Симметрия вторых оснований первого октета и центральной части второго октета}
\label{tab:degeneracy_symm}
\begin{tabular}{|l|l|c|c|c|c|c|c|c|c|}
\hline
 & продукт & G & A & S & P & V & T & L & R \\
\cline{2-10}
Октет \textbf{I}, группа \textbf{IV} & $M_{\text{side}}$ & 1 & 15 & 31 & 41 & 43 & 45 & 57 & 100 \\
\cline{2-10}
 & основание 2 & R & Y & Y & Y & Y & Y & Y & R \\
\hline
\hline
 & основание 2 & Y & R & R & R & R & R & R & Y \\
\cline{2-10}
Октет \textbf{II}, группа \textbf{II} & $M_{\text{side}}$ & 57 & 58 & 59 & 72 & 72 & 73 & 81 & 91 \\
\cline{2-10}
 & продукт & L & N & D & Q & K & E & H & F \\
\hline
\end{tabular}
\end{table}

\paragraph{Сигнатура 3.4B: Виртуальные кодоны.}

\begin{table}[htbp]
\centering
\caption{Идеограмма второго октета симметризованного генетического кода и виртуальные кодоны (пояснение в тексте)}
\label{tab:virtual_codons}
\footnotesize
\begin{tabular}{|c|c|c|c|c|c|c|c|c|c|c|c|c|c|c|c|c|c|}
\hline
группа & \multicolumn{2}{c|}{\textbf{III}} & \multicolumn{12}{c|}{\textbf{II}} & \multicolumn{2}{c|}{\textbf{I}} & $\Sigma$ \\
\hline
продукт & \textbf{C} & \textbf{I} & \textbf{стоп} & \textbf{S} & \textbf{L} & \textbf{N} & \textbf{D} & \textbf{Q} & \textbf{K} & \textbf{E} & \textbf{H} & \textbf{F} & \textbf{R} & \textbf{Y} & \textbf{M} & \textbf{W} & \\
$M_{\text{side}}$ & 47 & 57 & 0 & 31 & 57 & 58 & 59 & 72 & 72 & 73 & 81 & 91 & 100 & 107 & 75 & 130 & \\
\hline
основание 2 & G & T & A & G & T & A & A & A & A & A & A & T & G & A & T & G & \\
\hline
в-продукт & $-$ & \multicolumn{2}{c|}{\textbf{стоп}} & \multicolumn{2}{c|}{\textbf{стоп}} & \multicolumn{4}{c|}{\textbf{K}} & \multicolumn{4}{c|}{\textbf{старт/M}} & \multicolumn{3}{c|}{\textbf{старт/M}} & \\
$M_{\text{side}}$ & $-$ & \multicolumn{2}{c|}{0} & \multicolumn{2}{c|}{0} & \multicolumn{4}{c|}{72} & \multicolumn{4}{c|}{75} & \multicolumn{3}{c|}{75} & 222 \\
\hline
основание 2 & G & T & A & G & T & A & A & A & A & A & T & G & A & T & G & & \\
в-продукт & $-$ & \multicolumn{3}{c|}{\textbf{S}} & \multicolumn{4}{c|}{\textbf{K}} & \multicolumn{4}{c|}{\textbf{K}} & \multicolumn{2}{c|}{\textbf{C}} & \multicolumn{2}{c|}{$-$} & \\
$M_{\text{side}}$ & $-$ & \multicolumn{3}{c|}{31} & \multicolumn{4}{c|}{72} & \multicolumn{4}{c|}{72} & \multicolumn{2}{c|}{47} & \multicolumn{2}{c|}{$-$} & 222 \\
\hline
\end{tabular}
\end{table}


\section{Четвёртый информационный уровень. Звенья цепей}

На этом уровне боковые цепи аминокислот разделяются на подструктуры в последовательных позициях вдоль цепи, обозначаемые греческими буквами $\beta$, $\gamma$, $\delta$, $\varepsilon$, $\zeta$, \ldots\ для позиций звеньев цепи, начиная от $\alpha$-углерода.

\subsection*{Сигнатура 4.1. Уровень $\beta$ первого октета}

\begin{table}[htbp]
\centering
\caption{Массы звеньев уровня $\beta$ для октета~\textbf{I}}
\label{tab:level_beta_octet1}
\begin{tabular}{cc}
\begin{tabular}{|l|c|c|c|c|c|c|c|c|}
\hline
аминокислота & \textbf{G} & \textbf{A} & \textbf{S} & \textbf{P} & \textbf{V} & \textbf{T} & \textbf{L} & \textbf{R} \\
\hline
уровень $\beta$ & H & CH$_3$ & CH$_2$ & C$_2$H$_3$ & CH & CH & CH$_2$ & CH$_2$ \\
масса & 1 & 15 & 14 & 27 & 13 & 13 & 14 & 14 \\
\hline
$\Sigma$ & \multicolumn{8}{c|}{111} \\
\hline
\end{tabular}
&
\raisebox{-0.5\height}{\includegraphics[width=2cm]{cal_tab20_beta_octetI.png}}
\end{tabular}
\end{table}

Сумма масс уровня $\beta$ для первого октета равна SM-числу 111.

\subsection*{Сигнатура 4.2. Уровни $\beta$ и $\gamma$ второго октета}

\begin{table}[htbp]
\centering
\caption{Сумма масс уровней $\beta$ и $\gamma$ всего октета~\textbf{II}}
\label{tab:level_betagamma_octet2}
\footnotesize
\begin{tabular}{|l|c|c|c|c|c|c|c|c|}
\hline
аминокислота & \textbf{F} & \textbf{L} & \textbf{Y} & \textbf{стоп} & \textbf{C} & \textbf{стоп} & \textbf{W} & \textbf{H} \quad $\to$ \\
\hline
уровень $\beta$ & CH$_2$ & CH$_2$ & CH$_2$ & $-$ & CH$_2$ & $-$ & CH$_2$ & CH$_2$ \\
масса & 14 & 14 & 14 & 0 & 14 & 0 & 14 & 14 \quad $\to$ \\
\hline
уровень $\gamma$ & C & CH & C & $-$ & SH & $-$ & CH & C \\
масса & 12 & 13 & 12 & 0 & 33 & 0 & 13 & 12 \quad $\to$ \\
\hline
\end{tabular}

\vspace{0.3em}

\begin{tabular}{|l|c|c|c|c|c|c|c|c|c|}
\hline
аминокислота & \textbf{Q} & \textbf{I} & \textbf{M} & \textbf{N} & \textbf{K} & \textbf{S} & \textbf{R} & \textbf{D} & \textbf{E} \\
\hline
уровень $\beta$ & CH$_2$ & CH & CH$_2$ & CH$_2$ & CH$_2$ & CH$_2$ & CH$_2$ & CH$_2$ & CH$_2$ \\
масса & 14 & 13 & 14 & 14 & 14 & 14 & 14 & 14 & 14 \\
\hline
уровень $\gamma$ & CH$_2$ & C$_2$H$_5$ & CH$_2$ & C & CH$_2$ & OH & CH$_2$ & C & CH$_2$ \\
масса & 14 & 29 & 14 & 12 & 14 & 17 & 14 & 12 & 14 \\
\hline
\end{tabular}

\vspace{0.3em}

\begin{tabular}{|l|c|}
\hline
$\Sigma$ & 444 \\
\hline
\end{tabular}
\hfill
\raisebox{-0.5\height}{\includegraphics[width=2cm]{cal_tab21_betagamma_octetII.png}}
\end{table}

\subsection*{Сигнатура 4.3. Кодоны с первым пиримидином или первым пурином}

\begin{table}[htbp]
\centering
\caption{Кодоны с первым пиримидином, сумма масс уровней $\beta$ и $\gamma$. Массы уровней указаны в скобках}
\label{tab:betagamma_pyrimidine}
\footnotesize
\begin{tabular}{|l|c|c|c|c|c|c|c|}
\hline
аминокислота & TTY \textbf{F} & TTR \textbf{L} & TCN \textbf{S} & TAY \textbf{Y} & TAR \textbf{стоп} & TGY \textbf{C} & $\to$ \\
\hline
уровень $\beta$ & CH$_2$\,(14) & CH$_2$\,(14) & CH$_2$\,(14) & CH$_2$\,(14) & $-$\,(0) & CH$_2$\,(14) & $\to$ \\
\hline
уровень $\gamma$ & C\,(12) & CH\,(13) & OH\,(17) & C\,(12) & $-$\,(0) & SH\,(33) & $\to$ \\
\hline
\hline
аминокислота & TGA \textbf{стоп} & TGG \textbf{W} & CTN \textbf{L} & CCN \textbf{P} & CAY \textbf{H} & CAR \textbf{Q} & CGR \textbf{R} \\
\hline
уровень $\beta$ & $-$\,(0) & CH$_2$\,(14) & CH$_2$\,(14) & C$_2$H$_3$\,(27) & CH$_2$\,(14) & CH$_2$\,(14) & CH$_2$\,(14) \\
\hline
уровень $\gamma$ & $-$\,(0) & C\,(12) & CH\,(13) & CH$_2$\,(14) & C\,(12) & CH$_2$\,(14) & CH$_2$\,(14) \\
\hline
\hline
$\Sigma$ & \multicolumn{7}{c|}{333} \\
\hline
\end{tabular}
\hfill
\raisebox{-0.5\height}{\includegraphics[width=2cm]{cal_tab22_first_pyrimidine_bg.png}}
\end{table}

\begin{table}[htbp]
\centering
\caption{Кодоны с первым пурином, сумма масс уровней $\beta$, $\delta$ и $\zeta$. Массы уровней указаны в скобках}
\label{tab:betadeltazeta_purine}
\footnotesize
\begin{tabular}{|l|c|c|c|c|c|c|c|}
\hline
аминокислота & ATH \textbf{I} & ATG \textbf{M} & ACN \textbf{T} & AAY \textbf{N} & AAR \textbf{K} & AGY \textbf{S} & $\to$ \\
\hline
уровень $\beta$ & CH\,(13) & CH$_2$\,(14) & CH\,(13) & CH$_2$\,(14) & CH$_2$\,(14) & CH$_2$\,(14) & $\to$ \\
\hline
уровень $\delta$ & CH$_3$\,(15) & S\,(32) & $-$\,(0) & ONH$_2$\,(32) & CH$_2$\,(14) & $-$\,(0) & $\to$ \\
\hline
уровень $\zeta$ & $-$\,(0) & $-$\,(0) & $-$\,(0) & $-$\,(0) & NH$_2$\,(16) & $-$\,(0) & $\to$ \\
\hline
\hline
аминокислота & AGR \textbf{R} & GTN \textbf{V} & GCN \textbf{A} & GAY \textbf{D} & GAR \textbf{E} & GGN \textbf{G} & \\
\hline
уровень $\beta$ & CH$_2$\,(14) & CH\,(13) & CH$_3$\,(15) & CH$_2$\,(14) & CH$_2$\,(14) & H\,(1) & \\
\hline
уровень $\delta$ & CH$_2$\,(14) & $-$\,(0) & $-$\,(0) & O$_2$H\,(33) & C\,(12) & $-$\,(0) & \\
\hline
уровень $\zeta$ & C\,(12) & $-$\,(0) & $-$\,(0) & $-$\,(0) & $-$\,(0) & $-$\,(0) & \\
\hline
\hline
$\Sigma$ & \multicolumn{7}{c|}{333} \\
\hline
\end{tabular}
\hfill
\raisebox{-0.5\height}{\includegraphics[width=2cm]{cal_tab23_first_purine_bdz.png}}
\end{table}


\section{Пятый информационный уровень. Симметрии АРС}

Молекулы аминоацил-тРНК-синтетаз (АРС) реализуют генетический код, связывая аминокислоты с соответствующими тРНК. Существует два класса: АРС-1 (мономеры) и АРС-2 (димеры), каждый из которых содержит по 10 аминокислот, за исключением того, что лизин принадлежит обоим классам. Искусственно приписав лизин к АРС-2, мы получаем две равные группы по 10 аминокислот в каждой.

\begin{table}[htbp]
\centering
\caption{Два класса АРС, представленные соответствующими аминокислотами с первыми буквами кодонов (см. текст). Стрелки указывают на возрастание масс аминокислот}
\label{tab:ARS_classes}
\begin{tabular}{|c|c|}
\hline
\textbf{АРС-1 (мономеры)} & \textbf{АРС-2 (димеры)} \\
\hline
V, C, L, I, Q, E, M, R, Y, W & G, A, S, P, T, N, D, K, H, F \\
$\uparrow$ возрастание массы & $\uparrow$ возрастание массы \\
\hline
\end{tabular}
\end{table}

Симметрия первых оснований обнаруживается при сдвиге записей на один столбец. Нуклеотиды G и C образуют отдельные столбцы (обозначенные чёрным фоном в оригинале), тогда как A и T занимают другие столбцы. Расположение симметрично относительно вертикальной инверсии и циклического сдвига.

\begin{table}[htbp]
\centering
\caption{Нумерация аминокислот, соответствующая АРС-1 и АРС-2 (см. текст). Возрастание массы аминокислот каждого класса АРС показано интенсивностью серого цвета соответствующей ячейки}
\label{tab:ARS_numbering}
\begin{tabular}{|c|c|c|c|c|c|c|c|c|c|c|}
\hline
\textbf{V} & \textbf{C} & \textbf{L} & \textbf{I} & \textbf{Q} & \textbf{E} & \textbf{M} & \textbf{R} & \textbf{Y} & \textbf{W} & ARS-1 \\
\hline
1 & 2 & 3 & 4 & 5 & 6 & 7 & 8 & 9 & 10 & $+$ \\
\hline
\textbf{G} & \textbf{A} & \textbf{S} & \textbf{P} & \textbf{T} & \textbf{N} & \textbf{D} & \textbf{K} & \textbf{H} & \textbf{F} & ARS-2 \\
\hline
$-10$ & $-9$ & $-8$ & $-7$ & $-6$ & $-5$ & $-4$ & $-3$ & $-2$ & $-1$ & $-$ \\
\hline
\end{tabular}
\end{table}

\begin{table}[htbp]
\centering
\caption{Каллиграмма-B. Для каждого основания C, T, A, G (следующих в порядке возрастания веса, представленного уровнем серого) аминокислоты, соответствующие триплетам с данным первым основанием, выписаны в порядке возрастания веса и пронумерованы согласно их позиции в последовательности, соответствующей классам АРС-1 ($+$, чёрный текст) и АРС-2 ($-$, серый текст) (см. таблицу~25)}
\label{tab:calligram_B}
\begin{tabular}{|c|c|c|c|c|c|c|}
\hline
C & $-7$, \textbf{P} & $+3$, \textbf{L} & $+5$, \textbf{Q} & $-2$, \textbf{H} & $+8$, \textbf{R} & $(3, 2)$ \\
\hline
T & $-8$, \textbf{S} & $+2$, \textbf{C} & $-1$, \textbf{F} & $+9$, \textbf{Y} & $+10$, \textbf{W} & $(3, 2)$ \\
\hline
A & $-6$, \textbf{T} & $+4$, \textbf{I} & $-5$, \textbf{N} & $-3$, \textbf{K} & $+7$, \textbf{M} & $(2, 3)$ \\
\hline
G & $-10$, \textbf{G} & $-9$, \textbf{A} & $+1$, \textbf{V} & $-4$, \textbf{D} & $+6$, \textbf{E} & $(2, 3)$ \\
\hline
$\Sigma$ & $-31$ & 0 & 0 & 0 & $+31$ & \\
\hline
 & $(0, 4)$ & $(3, 1)$ & $(2, 2)$ & $(1, 3)$ & $(4, 0)$ & \\
\hline
\end{tabular}
\end{table}

Моделирование методом Монте-Карло даёт следующие вероятности для этих симметрий: $P_{\text{ас}} = 1{,}11\%$ (антисимметрия сумм столбцов), $P_{\text{стлб.сим}} = 18{,}8\%$ (столбцовая симметрия распределения типов АРС), $P_{\text{стр.сим}} = 34{,}4\%$ (строковая симметрия распределения типов АРС), с совместной вероятностью $P = 0{,}378\%$.


\section{Обсуждение}

\subsection{Осмысленные информационные структуры}

Представляют ли SM-сигнатуры осмысленную информацию? Все числа от 111 до 999 были обнаружены среди сигнатур. Теорема Пифагора $3^2 + 4^2 = 5^2$ появляется в первом октете. Для сравнения, сообщения Cosmic Call, отправленные внеземному разуму, также используют математические структуры в качестве сигналов привлечения внимания.

Информационные уровни 1--5 естественным образом возникли в нашем анализе, причём каждый уровень требовал всё более сложных конструкций и новых ресурсов. Ситуация напоминает головоломку~--- подобную судоку или системе диофантовых уравнений,~--- где обнаружение сигнатур на одном уровне мотивирует поиск на более глубоких уровнях.

\subsection{Об оценке степени <<невероятности>>}

Совокупная вероятность первых двух уровней составляет приблизительно $6 \times 10^{-9}$. 17 независимых SM-чисел на третьем уровне дают $(1/111)^{17} \approx 2 \times 10^{-35}$. Полная вероятность составляет приблизительно $10^{-45}$ или даже $10^{-50}$.

Однако <<проблема поиска по углам>> делает корректную оценку принципиально невозможной. Аналогия~--- обезьяна за пишущей машинкой: если мы ищем \emph{любую} удивительную закономерность, а не конкретную, мы ищем по углам обширного пространства возможностей, и вероятность найти \emph{что-нибудь} удивительное гораздо выше, чем вероятность любого конкретного удивительного результата.

Наш вывод состоит в том, что не существует объективного способа оценить невероятность наблюдаемых информационных структур, основываясь только на оценках вероятности.

\subsection{Перспективы решения проблемы случайность/искусственность}

Гипотеза искусственности фальсифицируема в смысле Поппера: если информационные сигнатуры могут быть объяснены биохимическими или эволюционными механизмами, гипотеза опровергнута. Гипотеза случайности, напротив, не фальсифицируема~--- любой набор числовых соотношений всегда может быть отвергнут как совпадение.

Из гипотезы искусственности следует конкретное предсказание: альтернативные генетические коды должны нарушать информационные сигнатуры. Мы исследовали 25 известных альтернативных генетических кодов. Из них 8 нарушают структуру Румера (первый информационный уровень), а остальные нарушают сигнатуры второго уровня. Это подтверждает первую половину предсказания. Вторая половина предсказания~--- что в альтернативных кодах не должны появляться альтернативные информационные сигнатуры~--- остаётся для проверки в будущем.


\section{Заключение}

В разделах~3--8 настоящей работы был представлен подробный обзор странных симметрий и информационных закономерностей генетического кода. Было обнаружено удивительно большое число сигнатур, и они естественным образом распределяются по пяти информационным уровням возрастающей сложности.

Первые два уровня особенно просты и играют роль сигналов привлечения внимания, аналогичных используемым в практике SETI. Информационные структуры не могут быть объяснены известными функциональными требованиями генетического кода.

Наиболее надёжным критерием установления искусственности было бы наличие семантического содержания, однако ситуация в этом отношении остаётся запутанной. Мы не можем полагаться только на оценки малых вероятностей из-за <<проблемы поиска по углам>>.

Гипотеза искусственности фальсифицируема через альтернативные генетические коды. Мы проверили 25 альтернативных кодов и подтвердили первую половину предсказания, вытекающего из гипотезы искусственности: все исследованные альтернативные коды нарушают информационные сигнатуры стандартного кода. Вопрос о том, обладают ли эти альтернативные коды собственными особыми информационными сигнатурами, остаётся открытым для будущих исследований.


\section*{Благодарности}

Авторы благодарят Максима Макукова за чрезвычайно плодотворное обсуждение.


\section*{Литература}

\begin{description}
\item Crick, F. (1968). The Origin of the Genetic Code. \textit{Journal of Molecular Biology}, 38, 367--379.
\item Danckwerts, H.J., and Neubert, D. (1975). Symmetries of Genetic Code-Doublets. \textit{Journal of Molecular Evolution}, 5, 327--332.
\item Dumas, S. (2007). The 1999 and 2003 Messages Explained. URL: \url{https://www.plover.com/misc/Dumas-Dutil/messages.pdf}.
\item Франк-Каменецкий, М. (2018). \textit{Самая главная молекула: От структуры ДНК к биомедицине XXI века}. Альпина Паблишер.
\item Hoesl, M., Oehm, S., Durkin, P., Darmon, E., Peil, L., Aerni, H., \textit{et al.} (2015). Chemical Evolution of a Bacterial Proteome. \textit{Angewandte Chemie}, 54, 10030--10034.
\item Hohsaka, T., Ashizuka, Y., Murakami, H., and Sisido, M. (2001). Five-Base Codons for Incorporation of Nonnatural Amino Acids into Proteins. \textit{Nucleic Acids Research}, 29, 3646--3651.
\item Конопельченко, Б., Румер, Ю. (1975). Классификация кодонов в генетическом коде. \textit{ДАН СССР}, 223, 471--474.
\item Makukov, M., and Shcherbak, V. (2018). SETI In Vivo: Testing the We-Are-Them Hypothesis. \textit{International Journal of Astrobiology}, 17, 1--20.
\item Marx, G. (1979). Message through Time. \textit{Acta Astronautica}, 6, 221--225.
\item McMurry, J. (2023). \textit{Organic Chemistry: A Tenth Edition}. OpenStax, Rice University, Houston, Texas.
\item Никитин, М. (2016). \textit{Происхождение жизни. От туманности до клетки}. М.: Альпина нон-фикшн.
\item Pacioli, L. (1508). \textit{De Viribus Quantitatis}. Library of the University of Bologna.
\item Румер, Ю. (2013). Систематизация кодонов в генетическом коде. \textit{ДАН СССР}, 183, 222--226.
\item Shcherbak, V., and Makukov, M. (2013). The `Wow! Signal' of the Terrestrial Genetic Code. \textit{Icarus}, 224, 228--242.
\item Wilhelm, T., and Nikolajewa, S.N. (2004). A New Classification Scheme of the Genetic Code. \textit{Journal of Molecular Evolution}, 59, 598--605.
\item Зайцев, А. (2008). Первое музыкальное межзвёздное радиосообщение. \textit{Радиотехника и электроника}, 53, 1107--1113.
\end{description}

\end{document}
